\begin{lemma}[Polynomial roots under specialization]\label{lem:hj-algebraic-root}
Let $Q\in\mathbb F[X,z,Y]$ be nonzero, with
$\deg_XQ\le T$, $\deg_zQ\le H$, and $1\le\deg_YQ\le B$.
Assume characteristic zero or characteristic $p>B$, and fix $D\ge0$.
Put
\[
 C=B+H,\qquad N=T+BD,\qquad
 I=\max\{0,2N-1\}C^2,\qquad
 J=C\bigl(1+\max\{0,2D-1\}C\bigr).
\]
Over $\overline{\mathbb F}$ there are at most $(2B+1)H$ exceptional
challenge values such that, outside them, all pairs $(z,P)$ satisfying
\[
 \deg P\le D,\qquad Q(X,z,P(X))\equiv0
\]
are contained in at most $I$ exceptional solution pairs and at most $C$
irreducible curves of actual polynomial solutions. In the affine space
of $z$ and the coefficients of $P$, these curves have total degree at
most $J$. Each retained curve is parametrized by an initial-value plane
curve through a map regular outside the exceptional challenge values.
\end{lemma}
\begin{proof}
Work over the algebraically closed constant field
$K=\overline{\mathbb F}$. Factor $Q$ in $K[X,z,Y]$, remove its
$Y$-independent content, and retain each positive-$Y$-degree irreducible
factor once. Call the resulting product $R$. Its separate degrees do
not increase. Since $p>B$ when the characteristic is positive, each
retained factor is separable in $Y$, so $R$ is squarefree over
$K(X,z)$. The removed content can specialize to zero identically in
$X$ at at most $H$ challenge values: exclude the roots of any one of
its nonzero $X$-coefficients. Outside these values, substitution into
$Q$ vanishes if and only if substitution into $R$ vanishes. If there
are no positive-$Y$-degree factors, this already proves the claim.

Choose a constant $\alpha\in K$ for which the $Y$-leading coefficient
of $R(\alpha,z,Y)$ and
$\operatorname{Res}_Y(R,R_Y)(\alpha,z)$ are nonzero polynomials in $z$.
Such a constant exists because only finitely many constants annihilate
all $z$-coefficients of either nonzero polynomial in $X,z$.
Their $z$-degrees are at most $H$ and $(2B-1)H$, respectively.
Exclude their roots as well. The combined exceptional set has size
at most $(2B+1)H$.

Write
\[
 F(z,u)=R(\alpha,z,u),\qquad A(z,u)=R_Y(\alpha,z,u).
\]
At every remaining point of $F=0$, one has $A\ne0$.
The nonvertical irreducible components of this plane curve have total
degree at most $C$, and there are at most $C$ of them. Vertical
components lie over the excluded leading-coefficient values.
On each nonvertical component the formal implicit solution has the
form
\[
 P(\alpha+t)=u+\sum_{j\ge1}a_jt^j,
 \qquad
 a_j=\frac{N_j(z,u)}{A(z,u)^{2j-1}},\qquad
 \deg N_j\le(2j-1)C.                         \tag{$*$}
\]
Here and below degrees in $(z,u)$ are total degrees.
For completeness, expand using Hasse coefficients,
$R(\alpha+t,z,u+v)=\sum c_{ab}(z,u)t^av^b$.
All $c_{ab}$ have degree at most $C$; on $F=0$ the constant term
vanishes and the linear term in $v$ is $Av$.
After solving the coefficient of $t^j$ for $a_j$, any remaining term
with $t$-degree $a$ and $v$-degree $b$ has denominator exponent
$e=2(j-a)-b+1\le2j-1$ and numerator degree at most $eC$,
by induction. Bringing it to denominator $A^{2j-1}$ proves $(*)$.
The term with $b=0$ has $a=j$ and denominator exponent one.
This recursion divides by $A$ only; it requires no factorial inverses.

Retain a component if its generic formal solution is a polynomial of
degree at most $D$. On every other component, some coefficient $a_j$
with $D<j\le N$ is nonzero. Indeed, otherwise the truncation $P_D$
through degree $D$ agrees with the formal solution through degree $N$.
The polynomial $R(\alpha+t,z,P_D(\alpha+t))$ then vanishes through
order $N$ and has degree at most $T+BD=N$, so it vanishes identically.
Uniqueness of the formal implicit root would make the original series
this degree-$D$ polynomial, a contradiction.
Choose one such $j$ on each discarded component. Every specialized
polynomial root of degree at most $D$ on it must satisfy $N_j=0$.
This numerator is not identically zero on that component.
B\'ezout's theorem, summed over components, bounds these initial
points by $\max\{0,2N-1\}C^2=I$.
Each regular initial point determines at most one polynomial solution,
so this counts solution pairs, not merely auxiliary points.

For $D\ge1$, put $\delta=A^{2D-1}$. On a retained component the map
to $(z,u,a_1,\ldots,a_D)$ has homogeneous coordinates represented by
\[
 \delta,\quad z\delta,\quad u\delta,\quad
 N_jA^{2D-2j}\quad(1\le j\le D).
\]
Their degrees are at most $1+(2D-1)C$. A rational map given by forms
of degree at most $L$ sends a plane curve of degree $d$ to a curve of
degree at most $dL$, by intersection with a generic hyperplane.
Summing proves the bound $J$; base points can only reduce it.
For $D=0$ the map is the identity on $(z,u)$ and the bound is $C$.
Changing from Taylor coefficients at $\alpha$ to coefficients in $X$
is an invertible linear transformation.
The maps are regular wherever $A\ne0$, and every such point of a
retained component is an actual polynomial solution: substitution of
the truncated polynomial is an identity in its function field, hence
at all points where its coefficients are regular. Their closures also
consist of actual solutions, since the substitution equations are
polynomial. Finally the maps remember both $z$ and $u$, and therefore
are generically injective. This proves all assertions.
\end{proof}
